% GENERATED FILE. Edit canonical lessons, then run npm run generate:books.
% canonical-source-hash: fnv1a64:2ac8e3cee62ab46b
% canonical-lessons: HI-C536-yojna, HI-C536-nirnay, HI-C536-vyavastha, HI-C536-ruprekha, HI-C536-masauda

\chapter{Making a Plan}
\label{ch:hi-making-a-plan}
\hlchaptermodality{536}

\begin{chapteropening}
\textbf{By the end of this chapter:} \emph{I can make a plan, sketch it out and say what is settled.}

Complete the last lesson of chapter 536: I can make a plan, sketch it out and say what is settled.
\end{chapteropening}

\section[\texorpdfstring{\hi{योजना} (yojnā)}{योजना (yojnā)}]{\hi{योजना} (yojnā) — a plan; a scheme}

\label{lesson:HI-C536-yojna}



\begin{quote}
Before the new one: say the Hindi for a medal, then the Hindi for to make run.
\end{quote}

\subsection*{You'll want to know: \hi{योजना}}
\textbf{\hi{योजना}} — \emph{yojnā} — \textquotedblleft{}a plan; a scheme\textquotedblright{}.

\begin{quote}
\hi{हमने} \hi{गर्मी} \hi{की} \hi{छुट्टी} \hi{के} \hi{लिए} \hi{एक} \hi{योजना} \hi{बनाई।}
\end{quote}

\begin{quote}
\emph{Hamne garmī kī chuṭṭī ke liye ek yojnā banāī.}
\end{quote}

\begin{quote}
We made a plan for the summer holiday.
\end{quote}

Feminine.

\begin{cousinweb}
\textbf{\hi{योजना}} (\emph{yojnā}) is a tatsama, taken whole from Sanskrit \emph{yojanā}, \textquotedblleft{}use, arrangement, preparation\textquotedblright{} (Monier-Williams). Hindi uses it for a plan.
\end{cousinweb}

\begin{grammarlens}[title={What you've built}]
One more word for saying what happens if.
\end{grammarlens}

\subsection*{Your turn}
\begin{itemize}
\raggedright
  \item \emph{Say it:} \emph{yojnā}
  \item \emph{Say it:} \emph{yojnā}, once more
  \item \emph{Say it:} \emph{Hamne garmī kī chuṭṭī ke liye ek yojnā banāī.} — We made a plan for the summer holiday.
  \item \emph{Recall:} say the Hindi for a medal, then the Hindi for to make run, then say \emph{yojnā} again
\end{itemize}

\begin{tcolorbox}[breakable,colback=teal!4,colframe=teal!35!black,title={Before you move on}]
What does \hi{योजना} mean? (\textquotedblleft{}A plan; a scheme\textquotedblright{}.) Say it once more.
\end{tcolorbox}
\section[\texorpdfstring{\hi{निर्णय} (nirṇay)}{निर्णय (nirṇay)}]{\hi{निर्णय} (nirṇay) — a decision (formal)}

\label{lesson:HI-C536-nirnay}



\begin{quote}
Before the new one: say the Hindi for to make run, then the Hindi for a plan.
\end{quote}

\subsection*{You'll want to know: \hi{निर्णय}}
\textbf{\hi{निर्णय}} — \emph{nirṇay} — \textquotedblleft{}a decision (formal)\textquotedblright{}.

\begin{quote}
\hi{हमने} \hi{गाँव} \hi{में} \hi{घर} \hi{बनाने} \hi{का} \hi{निर्णय} \hi{लिया।}
\end{quote}

\begin{quote}
\emph{Hamne gā̃v meṁ ghar banāne kā nirṇay liyā.}
\end{quote}

\begin{quote}
We decided to build a house in the village.
\end{quote}

Masculine.

A formal word; everyday speech says \hi{फ़ैसला}.

\begin{cousinweb}
\textbf{\hi{निर्णय}} (\emph{nirṇay}) is a tatsama, taken whole from Sanskrit \emph{nirṇaya}, \textquotedblleft{}decision, determination, settlement\textquotedblright{} (Monier-Williams).
\end{cousinweb}

\begin{grammarlens}[title={What you've built}]
One more word for saying what happens if.
\end{grammarlens}

\subsection*{Your turn}
\begin{itemize}
\raggedright
  \item \emph{Say it:} \emph{nirṇay}
  \item \emph{Say it:} \emph{nirṇay}, once more
  \item \emph{Say it:} \emph{Hamne gā̃v meṁ ghar banāne kā nirṇay liyā.} — We decided to build a house in the village.
  \item \emph{Recall:} say the Hindi for to make run, then the Hindi for a plan, then say \emph{nirṇay} again
\end{itemize}

\begin{tcolorbox}[breakable,colback=teal!4,colframe=teal!35!black,title={Before you move on}]
What does \hi{निर्णय} mean? (\textquotedblleft{}A decision (formal)\textquotedblright{}.) Say it once more.
\end{tcolorbox}
\section[\texorpdfstring{\hi{व्यवस्था} (vyavasthā)}{व्यवस्था (vyavasthā)}]{\hi{व्यवस्था} (vyavasthā) — an arrangement, a system}

\label{lesson:HI-C536-vyavastha}



\begin{quote}
Before the new one: say the Hindi for a plan, then the Hindi for a decision (formal).
\end{quote}

\subsection*{You'll want to know: \hi{व्यवस्था}}
\textbf{\hi{व्यवस्था}} — \emph{vyavasthā} — \textquotedblleft{}an arrangement, a system\textquotedblright{}.

\begin{quote}
\hi{शादी} \hi{में} \hi{खाने} \hi{की} \hi{व्यवस्था} \hi{बहुत} \hi{अच्छी} \hi{थी।}
\end{quote}

\begin{quote}
\emph{Śādī meṁ khāne kī vyavasthā bahut acchī thī.}
\end{quote}

\begin{quote}
The food arrangements at the wedding were very good.
\end{quote}

Feminine.

\begin{cousinweb}
\textbf{\hi{व्यवस्था}} (\emph{vyavasthā}) is a tatsama, taken whole from Sanskrit \emph{vyavasthā}, \textquotedblleft{}settlement, a fixed rule\textquotedblright{} (Monier-Williams).
\end{cousinweb}

\begin{grammarlens}[title={What you've built}]
One more word for saying what happens if.
\end{grammarlens}

\subsection*{Your turn}
\begin{itemize}
\raggedright
  \item \emph{Say it:} \emph{vyavasthā}
  \item \emph{Say it:} \emph{vyavasthā}, once more
  \item \emph{Say it:} \emph{Śādī meṁ khāne kī vyavasthā bahut acchī thī.} — The food arrangements at the wedding were very good.
  \item \emph{Recall:} say the Hindi for a plan, then the Hindi for a decision (formal), then say \emph{vyavasthā} again
\end{itemize}

\begin{tcolorbox}[breakable,colback=teal!4,colframe=teal!35!black,title={Before you move on}]
What does \hi{व्यवस्था} mean? (\textquotedblleft{}An arrangement, a system\textquotedblright{}.) Say it once more.
\end{tcolorbox}
\section[\texorpdfstring{\hi{रूपरेखा} (rūprekhā)}{रूपरेखा (rūprekhā)}]{\hi{रूपरेखा} (rūprekhā) — an outline; a rough plan}

\label{lesson:HI-C536-ruprekha}



\begin{quote}
Before the new one: say the Hindi for a decision (formal), then the Hindi for an arrangement, a system.
\end{quote}

\subsection*{You'll want to know: \hi{रूपरेखा}}
\textbf{\hi{रूपरेखा}} — \emph{rūprekhā} — \textquotedblleft{}an outline; a rough plan\textquotedblright{}.

\begin{quote}
\hi{अरुण} \hi{ने} \hi{योजना} \hi{की} \hi{एक} \hi{रूपरेखा} \hi{बनाई।}
\end{quote}

\begin{quote}
\emph{Aruṇ ne yojnā kī ek rūprekhā banāī.}
\end{quote}

\begin{quote}
Arun drew up an outline of the plan.
\end{quote}

Feminine.

\begin{cousinweb}
\textbf{\hi{रूपरेखा}} (\emph{rūprekhā}) is a modern compound of two Sanskrit words, \emph{rūpa}, \textquotedblleft{}outward appearance, form\textquotedblright{}, and \emph{rekhā}, \textquotedblleft{}a line\textquotedblright{} (Monier-Williams): the lines that give a plan its shape.
\end{cousinweb}

\begin{grammarlens}[title={What you've built}]
One more word for saying what happens if.
\end{grammarlens}

\subsection*{Your turn}
\begin{itemize}
\raggedright
  \item \emph{Say it:} \emph{rūprekhā}
  \item \emph{Say it:} \emph{rūprekhā}, once more
  \item \emph{Say it:} \emph{Aruṇ ne yojnā kī ek rūprekhā banāī.} — Arun drew up an outline of the plan.
  \item \emph{Recall:} say the Hindi for a decision (formal), then the Hindi for an arrangement, a system, then say \emph{rūprekhā} again
\end{itemize}

\begin{tcolorbox}[breakable,colback=teal!4,colframe=teal!35!black,title={Before you move on}]
What does \hi{रूपरेखा} mean? (\textquotedblleft{}An outline; a rough plan\textquotedblright{}.) Say it once more.
\end{tcolorbox}
\section[\texorpdfstring{\hi{मसौदा} (masaudā)}{मसौदा (masaudā)}]{\hi{मसौदा} (masaudā) — a draft (of a text)}

\label{lesson:HI-C536-masauda}



\begin{quote}
Before the new one: say the Hindi for an arrangement, a system, then the Hindi for an outline.
\end{quote}

\subsection*{You'll want to know: \hi{मसौदा}}
\textbf{\hi{मसौदा}} — \emph{masaudā} — \textquotedblleft{}a draft (of a text)\textquotedblright{}.

\begin{quote}
\hi{दफ़्तर} \hi{ने} \hi{नई} \hi{योजना} \hi{का} \hi{मसौदा} \hi{भेजा।}
\end{quote}

\begin{quote}
\emph{Daftar ne naī yojnā kā masaudā bhejā.}
\end{quote}

\begin{quote}
The office sent a draft of the new scheme.
\end{quote}

Masculine.

\begin{cousinweb}
\textbf{\hi{मसौदा}} (\emph{masaudā}) is usually traced to Arabic \emph{musawwada}, \textquotedblleft{}a rough copy\textquotedblright{}, through Persian.
\end{cousinweb}

\begin{grammarlens}[title={What you've built}]
One more word for saying what happens if.
\end{grammarlens}

\subsection*{Your turn}
\begin{itemize}
\raggedright
  \item \emph{Say it:} \emph{masaudā}
  \item \emph{Say it:} \emph{masaudā}, once more
  \item \emph{Say it:} \emph{Daftar ne naī yojnā kā masaudā bhejā.} — The office sent a draft of the new scheme.
  \item \emph{Recall:} say the Hindi for an arrangement, a system, then the Hindi for an outline, then say \emph{masaudā} again
\end{itemize}

\begin{tcolorbox}[breakable,colback=teal!4,colframe=teal!35!black,title={Before you move on}]
What does \hi{मसौदा} mean? (\textquotedblleft{}A draft (of a text)\textquotedblright{}.) Say it once more.
\end{tcolorbox}
